\chapter{Getting started with FHI-aims}
\label{Ch:quickstart}

\section{First step: Installation}

FHI-aims comes as a gzipped tar archive that can be extracted in any directory
of your choice, e.g., by typing

\begin{verbatim}
  gzip -d fhi-aims.tar.gz
  tar -xvf fhi-aims.tar
\end{verbatim}

at the
command line of any current Unix-like system.

\begin{center}
  \parbox[c]{0.8\textwidth}
  {\small
    Before you ask: FHI-aims is designed to run on any current Unix-based or
    Unix-like system, such as Linux, IBM's aix, Sun's Solaris, or even Mac OS
    X. However, we do \emph{not} support FHI-aims on Windows at this point. It
    is certainly possible to make it run on a Windows platform using the
    appropriate tools, but not simply out-of-the-box. 
  } 
\end{center}

The full package then extracts itself into a directory \emph{./fhi-aims}, with the
following subdirectories:
\begin{itemize}
  \item \emph{bin/} : Location for any FHI-aims binaries built using the
    standard Makefile
  \item \emph{doc/} : Contains possible further documentation.
  \item \emph{species\_defaults/} : Grids, basis sets and other defaults for
    chemical elements 1-102. These can be copy-pasted as ``species'' into the
    FHI-aims input file \texttt{control.in}. FHI-aims provides three levels of
    standard species defaults: ``light'', ``tight'',
    and ``really\_tight'' (see Sec. \ref{Sec:species}). In addition,
    some further preconstructed special-purpose species defaults are
    provided in a ``non-standard'' subdirectory.
  \item \emph{src/} : This directory, and its subdirectories, contain all FHI-aims source
    code files, as well as the Makefile needed to build this code. 
  \item \emph{testcases/} : Simple examples to test and illustrate the basic
    functioning of the code. The input files provided here may also be used as
    templates for any new electronic structure calculations, rather than
    assembling them from scratch.
  \item \emph{utilities/} : Some simple scripts to extract basic
    information from the standard output of FHI-aims:
    Visualization of geometries using the .xyz format, extracting a
    series of geometries during relaxation as a movie, or extracting
    the development of energies and forces during relaxation. There is
    also some more sophisticated infrastructure here: Script-based
    ab initio replica exchange molecular dynamics (Luca Ghiringhelli)
    and a basin-hopping framework to predict the structure of small
    clusters from scratch (Ralf Gehrke).
  \item \emph{regression/} : This directory contains a set of small standard
    test cases that can be run automatically using a script, \texttt{regression.py} --
    when run without any flags, it will provide its own self-documentation.
    We provide these regression tests purely for convenience. They
    will allow to check the compiled FHI-aims binary a little more
    extensively, but they are not strictly necessary to run
    FHI-aims. In particular, please do \emph{not} view the input files
    of the regression tests as FHI-aims best practices. Follow the
    manual, not simply the regression tests. In many cases,
    they are not. Rather, what is tested may be a corner case that can
    be handled differently (better) in normal practical scenarios.
\end{itemize}
A \texttt{README} file in that directory contains some of the quickstart information
given here in condensed format. 

\section{Building the code}

\emph{To state this at the outset: The ``aimsclub'' Wiki is the
  appropriate place to look for detailed compiler settings for
  specific platforms. Do not be shy to add your own, and do not be shy
  to ask in the Forum when in doubt. The information given in this chapter is
  essential as it explains the process, but the platform specific
  remarks in the Wiki may help you save some time.}

Since FHI-aims is distributed in source code form, the first task is to
compile an executable program. For this, the following \emph{mandatory prerequisites}
are needed:
\begin{itemize}
\item A working Fortran 95 (or later) compiler. A good example for x86 type
  computers is Intel's \texttt{ifort} compiler.  A free but significantly slower
  compiler for all platforms is \texttt{gfortran} from the GNU compiler
  collection (\url{http://gcc.gnu.org/fortran}) or the \texttt{g95}
  compiler (\url{http://www.g95.org}). Do not underestimate this
  slowdown, though -- a factor of three or so is possible.
  \item A compiled version of the \texttt{lapack} library, and a library
    providing optimized basic linear algebra subroutines (BLAS). Standard
    commercial libraries such as Intel's \texttt{mkl} or IBM's \texttt{essl}
    provide both \texttt{lapack} and BLAS support. \texttt{lapack} can also be found at
    \url{http://www.netlib.org/lapack/}. \\
    \emph{Having an optimized BLAS library for YOUR
    specific computer system(s) is critical for the performance of FHI-aims.}
    Very good free implementations include \texttt{ATLAS}
    (\url{http://math-atlas.sourceforge.net/}). 
\end{itemize}
You should also have a version of the Gnu \texttt{make} command, as not all
ancient Unix \texttt{make} commands support our \texttt{Makefile}. Typically, Gnu \texttt{make}
will already be present on your system, either as \texttt{make}, or possibly
as \texttt{gmake}. 

The next two prerequisites are \emph{optional}, but absolutely
encouraged for any advanced use of FHI-aims: Support for parallel
architectures, and (separately) support for fully parallel
Scalapack-based eigenvalue solvers. In these cases, you will also
need: 
\begin{itemize}
  \item A version of MPI libraries for parallel execution, often already
    present on a parallel system (if not,
    \url{http://www.open-mpi.org/} provides one of several free implementations)
  \item Compiled versions of the \texttt{scalapack} library, and
     basic linear algebra \emph{communication} subroutines (BLACS). Capable
     implementations can be found at \url{http://www.netlib.org/}, but
     are often provided already in the numerical libraries of many
     vendors (e.g., Intel MKL on Linux).
\end{itemize}
\textbf{The creation of a complete, MPI- and scalapack-enabled binary
  is effort well spent. This should be the goal when compiling FHI-aims
  for any production purposes. This means that you should ultimately
  aim to build the Makefile target (see below) called ``scalapack.mpi''.} 

The ``aimsclub'' Wiki gives compilation hints for many platforms
(including MPI, BLACS and Scalapack, which we strongly recommend on
any platform with more than one CPU). Please help by adding your own
working settings there, and ask us if there are questions.

Once all prerequisites are in place, change directory to the \emph{src/}
directory, and open the \texttt{Makefile} in a text editor. 

\begin{center}
  \parbox[c]{0.8\textwidth}
  {
  \emph{You must
  adjust at least some system-specific portions of the} \texttt{Makefile}\emph{---simply typing
  ``make'' and hoping that the problem will go away will not work.}
  }
\end{center}

Usually, all you will have to do is to decide on one of the preconfigured make
targets -- ``serial'', ``mpi'', or ``scalapack.mpi''. Near the top of the
\texttt{Makefile}, a number of \emph{mandatory} settings are commented for each
target. Uncomment \emph{only} the block of settings relevant to your chosen
make target, and fill in the correct values of each variable (\texttt{FC},
\texttt{FFLAGS}, \texttt{LAPACKBLAS}, ...) for your computer
system. \emph{Note} that the Makefile itself contains detailed instructions
and explanations regarding the meaning of these variables. Often, simply
adjusting the compiler name, the location of your libraries
(\texttt{LAPACKBLAS}, possibly \texttt{MPIFC} or \texttt{SCALAPACK}) will be
sufficient. In addition, we \emph{strongly} recommend that you consult the
documentation of  your compiler, in order to find out which optimization
options \emph{beyond} the generic ``-O3'' optimization level suggested preset in the
\texttt{Makefile} will make a  difference on your computer.

Finally, this brings us to the key step of the build process: Building the
code. After the \texttt{Makefile} is adjusted, type

\begin{verbatim}
  make <target>
\end{verbatim}

at the command line, where ``$\langle$target$\rangle$'' should be replaced by
the target of your choice: 
``serial'', ``mpi'', or ``scalapack.mpi''.\footnote{There is an
  additional target, \texttt{parser}, which builds an executable 
that stops after parsing the input. This binary can be used to check the
validity of input files. (The \keyword{dry\_run} keyword achieves the
almost same effect with the full binary.)
} 
If successful, this should build
the desired FHI-aims binary (the compilation will take a while) and place it
in the \emph{bin/} directory mentioned above.

Building FHI-aims can take a while nowadays. If you have more than one
processor on the machine for building FHI-aims, try

\begin{verbatim}
  make -j <number_of_processors> <target>
\end{verbatim}

This choice should speed up the process greatly.

You may also wish to keep your own copy of the \texttt{Makefile}, for
instance to be able to work directly with the FHI-aims git repository
without overwriting the general \texttt{Makefile}. In that case, just
copy the standard \texttt{Makefile} to something like
\texttt{Makefile.myname}, and use 

\begin{verbatim}
  make -j <number_of_processors> -f Makefile.myname <target>
\end{verbatim}

Finally, we do note that there is some support for more sophisticated
tasks, such as cross-platform builds (building binaries for different
architectures from the same home directory) among the non-standard
environment variables in the Makefile (see there). 

\section{Running FHI-aims}
\label{Sec:running}


As a simple test run to establish the correct functioning of FHI-aims and also
to familiarize yourself with the basic structure of the input and output
files, we suggest you change directories to the
\emph{testcases/H2O-relaxation/} directory. The test run provided there relaxes
a simple H$_2$O molecule from an initial (distorted) structure to the stable
one, and computes total energies, eigenvalues etc. along the way. Notice that
the key convergence settings (basis sets and grids) in this example are chosen
to be fast. The results (particularly the relaxed geometry) are still
trustworthy, but we encourage you already here to explore more
stringent convergence settings later. In fact, \emph{always} explore the
impact of critical convergence settings on the accuracy of key results in your
own project.

In the \emph{testcases/H2O-relaxation/} directory, type

{
  \verb+ ../../bin/aims.+\emph{version}\verb+.serial.x < /dev/null | tee H2O_test.own+
}

at the command line. For ``\emph{version}'', you must insert the code version stamp 
that was actually downloaded and built (for example, \verb+010110+, \verb+071711_5+, \verb+081912+, 
\verb+081213+, \verb+071914+, or whichever code version you are building).\footnote{The 
FHI-aims version stamp can be modified to whatever you wish in \texttt{Makefile.backend} in the \emph{src/} directory.}
Similarly, for make targets other than ``serial'', you should
use the appropriate binary including the necessary mpi command
instead. On many (but not all) platforms, that command will be
\texttt{mpirun}, and will also require you to specify the number of processors
to be used by a flag. 

The result will be an output stream on your computer screen (created by
``tee'') which is also captured in an output file \texttt{H2O\_test.own}. 
Any critical information regarding computational
settings, results (total energies, forces, geometries, ...), errors
etc. should be contained in this file, which we encourage you to look at (yes,
it is \emph{meant} to contain human-readable and useful explanations). Any \emph{other} output files are  
only written if requested, and will be covered in the later sections of this
text.

\begin{center}
  \parbox[c]{0.8\textwidth}
  {
  \emph{The standard output stream or file contains any and all output that FHI-aims
  writes by default. For later use, you must save this output stream to disk in some way, 
  using standard Unix redirections such as the \texttt{tee} command above or a simple 
  redirect.}
  }
\end{center}

  Apart from the first expression given above, such 
  redirections might look like this: 

  {
    \verb+ ../../bin/aims.+\emph{version}\verb+.serial.x < /dev/null > H2O_test.own+
  }

  or even like this: 

  {
    \verb+ nohup ../../bin/aims.+\emph{version}\verb+.serial.x < /dev/null > H2O_test.own 2>&1 &+
  }

  The latter version decouples the FHI-aims run completely from your current login, and additionally
  saves any system error messages to the standard output file as well. With the above command sequence,
  you may safely log out from the computer in question, the code should keep running in the background. 
  Take care to monitor your running processes using the \texttt{ps} Unix command. For instance, it is
  highly unadvisable to run ten instances of FHI-aims at once in the background on a single CPU and expect
  any reasonable performance of the computer at all. The above hints
  are just examples of general Unix command-line sequences. For a complete treatment, we
  recommend that beginners read a separate Unix textbook, or---often feasible---learn by doing and Google. 

If successful (otherwise, consider the warnings three paragraphs below), you
may wish to compare your results to those contained in our 
own output from this run, which is contained in the file
\texttt{H2O.reference.out}. You should obtain exactly the same total energies,
forces, and geometries as given in this file. Any information regarding timing
is, of course, specific to your computer environment, and not necessarily the
same. 

The directory \emph{testcases/H2O-relaxation/} contains two more files,
\texttt{control.in} and \linebreak[4] \texttt{geometry.in}. These are the sole two input
files required by FHI-aims, and are the most important files to learn about in
the rest of this documentation. In brief, \texttt{geometry.in} contains any
information related directly to a system's geometry -- normally, this will be
atomic positions (in \AA) and perhaps lattice vectors for periodic
calculations, but no more. Any other, method-related, input information is
part of \texttt{control.in}. 

In practice, we attempt to strike a balance
between the information \emph{needed} by \linebreak[4] \texttt{control.in}, and information
set to safe defaults unless specified explicitly. For example, you \emph{must}
specify the level of theory (e.g., the fact that PBE exchange-correlation is
used) and also the basis set employed. While it is highly useful to
have this relevant information openly accessible, this  
would also create the need to personally edit a large amount of input before ever
tackling the first run. For any information tied to the actual element (or
``species''; arguably the most complex information required), we therefore provide
ready-made template files for all elements (1-102) in the
\emph{species\_defaults} directory. They are ready for copy-paste into
\texttt{control.in}. These files will still benefit from some
adjustment to your personal needs (for instance, the provided integration
grids are set rather on the safe side, at the expense of more CPU time), but
should greatly simplify the task. 

Two final, important warnings regarding the execution of FHI-aims that are beyond
our direct control:
\begin{itemize}
  \item FHI-aims \emph{requires} that the execution stack size available to
    you be large enough for some initial internal operations. Spare us the
    details (ample explanation of the meaning of the ``stack'' in Unix can be
    found elsewhere), but for reasons unbeknownst to us, some vendors limit
    the default user stack size to $\approx$5~MB in a time when typical
    available RAM per processor is 2 GB or more. If too little stack is
    available, your FHI-aims run will \emph{segfault} shortly after the
    command was launched. In that case, type:
    \begin{verbatim}
        ulimit -s unlimited
    \end{verbatim}
    (when using the bash shell or similar), or
    \begin{verbatim}
        limit stacksize unlimited
    \end{verbatim}
    (when using the tcsh or similar). 
    \begin{verbatim}
        echo $SHELL
    \end{verbatim}
    will tell you which shell you are using. Ideally, this same setting should
    be specified in your .profile, .bashrc, or .cshrc login profiles. If
    ``unlimited'' does not work, try setting a large value instead, e.g.,
    \texttt{ulimit -s 500000}.
  \item An important system settings for parallel execution is the environment
    variable  
    \begin{verbatim}
      export OMP_NUM_THREADS=1
    \end{verbatim}
    (the syntax is correct for the bash shell). When using Intel's mkl, you
    should additionally set \texttt{MKL\_NUM\_THREADS} to 1 and
    \texttt{MKL\_DYNAMIC} to FALSE.
\end{itemize}

After startup, the first messages contain information about your computer's 
environment: Code version, compiler information, host names, environment 
variables which turned out to be useful and which should be set on
your system (e.g.~\texttt{OMP\_NUM\_THREADS}), etc. The complete input
files \texttt{control.in} and \texttt{geometry.in} are also repeated
verbatim. Any FHI-aims run should thus be completely reproducible
based on the standard output stream alone.

Should you encounter further issues, consider also the troubleshooting
information documented in Appendix \ref{appendix_trouble_shooting}. 

All this said, after successfully running the test run, you should now be
ready to go with FHI-aims. The remainder of this document is about the details
-- available options, how to run aims most efficiently, etc. Happy computing!


\section{Shortcut: make.sys}

The initial set of builds described so far worked by directly editing
the \texttt{Makefile} shipped with FHI-aims. This is useful because
working examples of Makefile variables are included there, as well as
the comments documenting the meaning of individual variables.

Once you are familiar with that process, there is a workaround that
does not require to edit the Makefile itself at all. Instead,
FHI-aims also allows a user to create a file \texttt{make.sys} in the
\emph{src} directory, containing only the Makefile variables that need
to be set. For example, a \texttt{make.sys} file that builds a
Fortran-only scalapack.mpi build could look like this:

\begin{verbatim}
 FC = ifort
 FFLAGS = -O3 -ip -module $(MODDIR)
 F90MINFLAGS = -O0 -module $(MODDIR)
 FMINFLAGS = -O0 -module $(MODDIR)
 KNOWN_COMPILER_BUGS = ifort-14.0.1-O3
 F90FLAGS = $(FFLAGS)
 ARCHITECTURE = Generic
 LAPACKBLAS = -L/opt/intel/composerxe/mkl/lib/intel64 \
   -I/opt/intel/Compiler/11.0/084/mkl/include/intel64/lp64 -lmkl_intel_lp64 \
   -lmkl_sequential -lmkl_core -lpthread
 USE_MPI = yes
 MPIFC = mpiifort
 SCALAPACK = -L/opt/intel/composerxe/mkl/lib/intel64 \
   -I/opt/intel/Compiler/11.0/084/mkl/include/intel64/lp64 -lmkl_scalapack_lp64 \
   -lmkl_blacs_intelmpi_lp64
\end{verbatim}

Note that this example of a \texttt{make.sys} file comes from a
specific computer and holds only for a specific compiler
version. Please do not just use it blindly -- instead, please
determine the compiler, libraries, compiler flags etc. that are
appropriate for your computer system when creating your own file
\texttt{make.sys}. 

\section{Cross-Compiling with a C Compiler}
\label{cross-compile-c}

The standard build process of FHI-aims is intended to work with
\emph{only} a Fortran compiler and a minimal number of libraries,
e.g., those needed towards the scalapack.mpi build target (see
above). The idea is that this provides all the necessary
infrastructure to run FHI-aims with optimal performance on a given
platform, yet does not overburden a user with the need to balance an
overwhelming set of potentially conflicting system or environment
variables if they are not needed.

That said, FHI-aims does provide some functionality that can only be
accessed by compiling part of the code base with a C compiler, in some
cases simply because a system call in question is not available from a
Fortran interface. 

Cross-compiling with a C compiler can be simple \emph{if} your Fortran
and C compiler use compatible interfaces. (This may not be always the
case.) In the simplest case, adding the following two variables to
your \texttt{Makefile} or to your \texttt{make.sys} file will enable
some basic functionality that can only be accessed with a C
cross-compiled build:

\begin{verbatim}
 USE_C_FILES = yes
 CC = gcc
\end{verbatim}

The CC variable should specify a C compiler compatible with your
Fortran compiler on your specific computer system. This choice may not
always be the \emph{gcc} compiler. Again, C cross-compiling is not
necessary, but if you do, please make sure that the right C compiler
is chosen for your system.

An additional variable \texttt{CCFLAGS} can be set to specify C
compiler flags along with the \texttt{CC} variable above and will add
C compiler flags that might be necessary for your particular computer setup.


\section{Compilation options beyond the standard Makefile}
\label{Sec:Makefiles}

The build provided by the standard \texttt{Makefile} in FHI-aims is designed
for minimal complexity to obtain the full functionality that most users should
have. Separate ``basic linear algebra subroutines'' (BLACS), Lapack,
parallel builds (if a parallel machine is available, nowadays almost always),
and scalapack support are so performance-critical that every user should spend
the time to investigate them in detail before doing serious production work
with FHI-aims. These dependencies of FHI-aims on external libraries are
therefore kept in the main \texttt{Makefile}.

In addition, FHI-aims provides further functionality that can be achieved by
linking to other external libraries. However, this functionality will not be
needed by all users and/or could seriously complicate the build process for
everyone. Such functionality is therefore available through separate, amended
versions of the \texttt{Makefile}. We encourage everyone to try these builds
(they are not so difficult after all), but they should not become stumbling
blocks. 

In particular, several optional Makefile with additional functionality
exist (and could even be combined):
\begin{itemize}
  \item \texttt{Makefile.cuba} : Allows to compile in the separate ``CUBA''
    Monte Carlo integration library, which enables the Langreth-Lundqvist van
    der Waals functional based on noloco as a post-processing step. See Section
    \ref{Sec:vdwdf} for more details.
  \item \texttt{Makefile.meta} : \emph{still experimental!} Allows to
    interface FHI-aims to the \texttt{PLUMED} library for free-energy
    calculations for molecular dynamics. \\ (see
    \url{http://merlino.mi.infn.it/\~plumed/PLUMED/Home.html} , the \texttt{PLUMED}
    project homepage). Currently, a copy of the \texttt{PLUMED} library is
    kept in the \emph{external} directory of the FHI-aims source code, and
    must be compiled separately using a C compiler. Note that
     the token ``-DFHIAIMS'' must be included in the CCFLAGS (at present, this is already specified in \texttt{plumed.inc}). 
     For compiling on IBM power machines, the flags -mpowerpc64 -maix64 should
    be included. In the future, this will be corrected by housing the
    respective FHI-aims plugin directly in the \texttt{PLUMED} library. We
    apologize that the linking process is not yet further documented,
    but if this functionality is of interest to you, please contact us.\\
    By selecting the correct target in \texttt{Makefile.meta}, the code can be compiled 
    with \texttt{lapack} or \texttt{scalapack} libraries, or with shared memory support (see next item).
  \item \texttt{Makefile.ipi} : Allows to interface with the i-PI path integral molecular dynamics wrapper\cite{CeriottiMoreMano_2013}.
    See the \keyword{use\_pimd\_wrapper} keyword for a few more details.
  \item \texttt{Makefile.shm} : Another example of a Makefile that
    cross-links C and Fortran based functionality, although the
    actual functionality that Makefile.shm provides --  access to
    shared-memory arrays in the Hartree potential -- is no longer
    needed at this point.
  \item \texttt{Makefile.amd64\_SSE} : This Makefile contains an
    example how to use FHI-aims together with a version of the 
    ELPA library which is especially tuned for SSE vector instructions.
    Note that the GNU C compiler must be installed in order to produce
    running code from the specific assembler files. As the name
    ``amd64'' indicates this optimization is only reasonable on
    processors which support the amd64 instruction set.
    Note that ''ELPA\_ST'' is set, to use the single-threaded version
    of ELPA.
  \item \texttt{Makefile.amd64\_SSE\_mt} : The same as
    Makefile.amd64\_SSE, but by setting ''ELPA\_MT'' the hybrid MPI
    OpenMP version of ELPA is used.
  \item \texttt{Makefile.amd64\_AVX} : This Makefile contains an
    example how to use FHI-aims together with a version of the 
    ELPA library which is especially tuned for AVX vector instructions.
    Note that the GNU C and C++ compiler must be installed in order to produce
    running code from the specific ELPA files which contain gcc
    intrinsic functions. As the name  ``amd64'' indicates this
    optimization is only reasonable on  processors which support the 
    amd64 instruction set and which already ''understand'' AVX,
    i.e. Intel Sandybridge or newer. Note that ''ELPA\_ST'' is set, to 
    use the single-threaded version of ELPA.
  \item \texttt{Makefile.amd64\_AVX\_mt} : The same as
    Makefile.amd64\_AVX, but by setting ''ELPA\_MT'' the hybrid MPI
    OpenMP version of ELPA is used.

\end{itemize}

\subsection{Alternative build process with the Waf build system}

\emph{still experimental!}

FHI-aims can be alternatively built with the Waf build system 
(\url{https://waf.io/}).
Waf is a modern build system written in Python which enables writing build 
scripts (also in Python) that are easier to maintain compared to makefiles.
Another advantage specific to Fortran is that the dependencies between modules 
are recognized automatically by Waf.
The actual build flow (such as recompiling only changed and depending files) and 
resulting binaries should be bit-for-bit identical to the standard build process 
with \verb+make+.

The Waf executable (a Python script) is distributed together with FHI-aims and 
is located at \verb+src/waf+.
The Waf build script is at \verb+src/wscript+ and is not supposed to be changed 
under normal circumstances.
The build script assumes that the file \verb+src/Makefile+ is unchanged from the 
distributed version and all user build environment variables (such as \verb+FC+ 
or \verb+LAPACKBLAS+) are defined in \verb+src/make.sys+ (which is sourced by 
the standard makefiles).
If this prerequisite is satisfied, building with Waf should be as easy as
\begin{verbatim}
cd src
./waf configure <options>
./waf install
\end{verbatim}
where \verb+<options>+ can be any of the AIMS build extensions such as 
\verb+--with-mpi+ or \verb+--with-scalapack+.
For a full list of available extensions, see
\verb+./waf --help+.
Both Waf calls can be joined to
\verb+./waf configure install <options>+.
The build is done in directory \verb+build+ and installs the binaries into 
\verb+bin+, identically to the standard build process.
Subsequent builds can be run only with
\begin{verbatim}
./waf install
\end{verbatim}

With the Waf build system, all build variants of AIMS (such as MPI, Scalapack, 
building AIMS as a shared library, any additional libraries) are managed in 
so-called build extensions defined in \verb+wext_<name>.py+ files.
A build extension defines changes to the basic build process both when it is 
activated (via the command-line option) and when it is not activated (there is 
often a need for so-called stub files).
To learn how to write an AIMS build extensions, see the \verb+wext+ module 
documentation below.

\paragraph{\texttt{wext} --- AIMS build extensions} \ 


The \hyperref[module-wext]{\texttt{wext}} module manages AIMS build
extensions such as MPI or Scalapack.

To make a new build extension, create a file called
\texttt{wext\_\textless{}name\textgreater{}.py} in the \texttt{src}
directory where \texttt{\textless{}name\textgreater{}} is the name of
your extension. This Python file can be empty, in which case the
extension has no effect besides adding an option to the Waf options (see
\texttt{./waf\ -\/-help}). The extension is then activated by calling
Waf with this option. Functionality is added by defining variables and
functions (all optional) in
\texttt{wext\_\textless{}name\textgreater{}.py}:

\begin{description}
\item[help]
A string with a help for the build extension which appears in
\texttt{./waf\ -\/-help}.
\item[require]
A list of build extensions that the extension requires if activated.
\item[stub(ctx)]
Function that takes a \hyperref[wext.Context]{\texttt{Context}} (see
below) as an argument and applies actions (such as adding source files)
that should happen if the extension is not active.
\item[extension(ctx)]
Takes a \hyperref[wext.Context]{\texttt{Context}} as argument and
applies actions that should happen when the extension is active.
\end{description}

For example, the Scalapack extension is defined as:

\begin{verbatim}
help = 'compile aims with Scalapack'
require = ['mpi']

def stub(ctx):
    ctx.add_files('fc', 'scalapack_stub.f90')

def extension(ctx):
    ctx.add_link_flag(ctx.env.SCALAPACK)
\end{verbatim}

In common language, Scalapack requires MPI to work. The Scalapack stub
adds an appropriate stub file (compiled with Fortran compiler -- `fc').
And finally if Scalapack is activated, a link flag is added from the
\texttt{SCALAPACK} environment variable. For full list of avaialble
actions, see documentation of \hyperref[wext.Context]{\texttt{Context}}.

\begin{description}
\item[ \emph{class} \texttt{wext.}\texttt{Context}{(}\emph{wctx}{)}]
Provides interface to modify a build process.

\begin{description}
\item[ \texttt{add\_feature}{(}\emph{feat}, \emph{files}{)}]
Add a feature to files.
\end{description}

\begin{description}
\item[ \texttt{add\_files}{(}\emph{feat}, \emph{files}{)}]
Add files with features to the build process.
\end{description}

\begin{description}
\item[ \texttt{add\_link\_flag}{(}\emph{flag}{)}]
Add a link flag.
\end{description}

\begin{description}
\item[ \texttt{change\_ext}{(}\emph{fext}{)}]
Change file name extension (including the dot) of the resulting binary.
\end{description}

\begin{description}
\item[ \texttt{env} \emph{= None}]
Provides access to the variable environment.
\end{description}

\begin{description}
\item[ \texttt{replace\_files}{(}\emph{newfiles}, \emph{oldfiles},
\emph{feat=None}{)}]
Replace old files with new files.

Optionally add a feature. All old files have to have the same features.
\end{description}

\begin{description}
\item[ \texttt{wctx} \emph{= None}]
Provides access to the Waf build configuration context.
\end{description}
\end{description}



All inquiries, issues and general questions about the Waf build system should be 
directed at 
\href{mailto:jhermann@fhi-berlin.mpg.de?subject=\%5Baims\%20waf\%5D}{jhermann@fhi-berlin.mpg.de}.

\section{Compiling faster versions of FHI-aims on specific platforms}
\label{Sec:Platform}

FHI-aims is intended to be a Fortran-only code, which -- for most of
the code -- means that building the ``fastest'' version of FHI-aims on
a given computer architecture is ``only'' a matter of finding the
right Fortran compiler and compiler options for that processor. For
some architectures, specific compiler options are collected in the
FHI-aims club wiki -- please check there and please add any useful
information that you may find.

That said, one particular performance-critical area for large systems
is the Kohn-Sham eigenvalue solver. In FHI-aims and on parallel
computers, this problem is solved by the ELPA library. ELPA, in fact,
allows its users to specify specific, platform-optimized so-called
linear algebra ``kernels.'' By default, FHI-aims uses a generic kernel
which will compile with any Fortran compiler and will give reasonable
speed. However, if one knows which specific computer chip one is
using, it is possible to substitute this kernel with an architecture
specific kernel and compile a faster version of ELPA into
FHI-aims. This is possible, for example, for the BlueGene/P,
BlueGene/Q, Intel AVX and several other Intel architectures. For
standard Intel x86 chips, there is even an ``assembler'' based kernel
that will get fast performance regardless of the Fortran compiler
above.

Note that this choice can matter. For example, the ``generic'' ELPA
kernel will produce fast code for the Intel Fortran compiler, but much
slower code with certain versions of the PGI Fortran compiler (often
found on Cray machines).

How to build FHI-aims with these special ELPA kernels?

In some cases, the necessary information already exists in the form of
architecture-specific Makefiles in FHI-aims, e.g.,
\texttt{Makefile.amd64\_SSE} or \texttt{Makefile.amd64\_AVX} mentioned
in the last section. In other cases (e.g., for the Intel Assembler
kernel elpa2\_kernels\_asm\_x86\_64.s, found in the src directory) you may
have to precompile with the Gnu assembler ``as'' before building the
rest of FHI-aims.

Going away from the ``generic'' kernels can therefore require a little
bit of extra work. The point of the present section is to clarify that
this extra work is possible, and often worthwhile, if you already know
that you are going to use FHI-aims for a long time on the same
architecture. 
